% !TEX root = main.tex
\section{Giới thiệu}
Các mô hình ngôn ngữ lớn (LLM) hiện đại đạt năng lực giải quyết bài toán phức tạp nhờ kết hợp hai yếu tố: kiến trúc Transformer tối ưu cho suy diễn tự hồi quy và khả năng phân bổ thêm tính toán tại thời điểm suy luận \cite{wei2022chain,wang2023selfconsistency,snell2024scaling}. Tuy nhiên, phần lớn các hệ thống mã nguồn mở hiện nay phụ thuộc vào các lớp đóng gói sẵn như \texttt{AutoModelForCausalLM} của Hugging Face, khiến người học khó nắm vững bản chất toán học của cơ chế chú ý, quản lý bộ nhớ KV Cache cũng như quy trình huấn luyện suy luận.

Trong đồ án này, chúng tôi tự lập trình toàn bộ hệ thống từ đầu bằng PyTorch thuần (\texttt{torch.nn.Module}) nhằm đáp ứng hai mục tiêu chính. Thứ nhất, ở tầng Khung Cốt lõi, hệ thống xây dựng bộ tách từ Byte-Level BPE tối ưu cho chữ số, kiến trúc Transformer hiện đại (~43,8M tham số chạy vừa trên Colab T4 16GB), và \texttt{DynamicKVCache} kèm kiểm chứng tương đương số học 2 cấp độ cũng như đo đạc profiling định lượng giữa chế độ bật và tắt KV Cache tại các độ dài ngữ cảnh $L \in \{128, 512, 1024, 2048\}$. Thứ hai, ở Hướng C về suy luận và tính toán thời điểm suy luận, mô hình được huấn luyện CoT SFT trên dữ liệu đã lọc rò rỉ theo định dạng \texttt{<think>...</think>} và \texttt{\#\#\#\# <answer>}, sau đó đánh giá đối chiếu có hệ thống giữa \texttt{Greedy Pass@1}, \texttt{Majority Voting (Maj@K)} và \texttt{Best-of-N (BoN@K)} trên cùng một nhóm $K=16$ chuỗi sinh thử trên toàn bộ 1.319 câu GSM8K Test \cite{cobbe2021gsm8k}.
